\section{问题重述}
\subsection{问题背景}
洪涝灾害可能造成道路中断，使常规运输难以在首批救援窗口内到达偏远村落。无人机可绕开地面通行障碍，但其载荷、续航、电池周转和通信距离均有限；山区地形起伏还会引起视距遮挡。因此，方案设计需要同时考虑物资批次、运输时序、通信保障及多组独立执行时的资源配置。

\subsection{问题内容}
题目给定应急物资货箱、服务区、调度中心、运输无人机及电池参数，并给出地形和通信数据。本文按题意将任务拆为四个相互衔接的问题：
\begin{enumerate}
  \item 对各服务区计算不同机型的安全载荷，并将不可拆货箱组成满足载荷、容积和能源约束的单区运输批次。
  \item 在运输批次基础上确定多区路线、起飞时间、无人机及电池分配，使需求获得服务并满足硬时限与资源周转条件。
  \item 对运输过程中的地空链路进行通信模式配置；在直连不足时引入空中中继，并统计中继任务与额外能源需求。
  \item 冻结既有运输和通信任务，将服务区分为 $K=2$ 或 $K=3$ 个独立任务组，保证同一跨区运输任务不被拆分，比较各组峰值并发资源与库存缺口。
\end{enumerate}

\subsection{技术路线}
本文采用“安全载荷与组批—时限排程—通信模式选择—固定方案分区核算”的递进式路线。前一问题输出作为后一问题的固定输入，使每问新增决策边界清晰。全流程示意见图~\ref{fig:route}。
\begin{figure}[htbp]
  \centering\includegraphics[width=0.96\linewidth]{technical_route.png}
  \caption{四问协同技术路线}\label{fig:route}
\end{figure}
